\documentclass[11pt,a4paper]{article}
\usepackage[utf8]{inputenc}
\usepackage[T1]{fontenc}
\usepackage[brazil]{babel}
\usepackage[a4paper,margin=2.3cm]{geometry}
\usepackage{enumitem}
\usepackage{hyperref}
\usepackage{xcolor}
\setlength{\parindent}{0pt}
\setlength{\parskip}{5pt}
\sloppy
\setlength{\emergencystretch}{3em}
\hypersetup{colorlinks=true,urlcolor=blue,linkcolor=black}
\newcommand{\cabecalho}[2]{\begin{center}{\small \textbf{Sprint 2 -- #1}}\\[1mm]\textit{Inventário de Ativos e Vulnerabilidades}\\Matheus Moreli Durão Silva -- 12621CBS217\\Cibersegurança -- Bacharelado -- Noturno\end{center}\hrule\vspace{4mm}}
\newcommand{\bloco}[1]{\par\vspace{2mm}\textbf{#1}\par}
\newcommand{\codigo}[1]{\texttt{#1}}

\begin{document}\cabecalho{S2\_16 -- Branches}{}
\bloco{O que é} Uma branch é uma linha de trabalho independente no mesmo repositório. A \codigo{main} deve representar uma versão estável; uma branch como \codigo{feature/s2-08-cadastro-ativos} concentra uma mudança específica sem afetar diretamente a principal.
\bloco{Fluxo sugerido por card} Atualizar a cópia local da \codigo{main}; criar branch com nome do card; fazer alteração coerente; testar; criar commit explicando a mudança; enviar a branch; abrir Pull Request; revisar; integrar e testar a \codigo{main}. Um card que só produz documentação também pode ter branch própria, por exemplo \codigo{docs/s2-04-requisitos-crud}.
\bloco{Importante} Não fazer commit diretamente na \codigo{main}. Antes de integrar, sincronizar a branch com a principal para perceber conflitos cedo. Um commit deve ter uma ideia verificável, não várias mudanças sem relação.
\bloco{Evidência esperada} Branch visível no GitHub, histórico de commits coerentes, Pull Request e teste após merge. A criação e o merge dependem da execução no repositório remoto.
\bloco{Checklist} Sem commit direto na main, nome representativo, commits coerentes, sincronização e integração verificada.
\end{document}
